\section{Introduction}
\label{sec:intro}

The proliferation of hyper-realistic facial manipulation algorithms has destabilized the evidential integrity of digital video media. Modern generative pipelines—spanning autoencoder-based face swaps \citep{faceswap2018}, neural texture synthesis \citep{thies2019neural}, and diffusion-based identity reenactment \citep{ho2020denoising, ramesh2022hierarchical}—now produce synthetic artifacts that are imperceptible to human scrutiny. Consequently, automated facial forgery detection has emerged as a critical societal safeguard for digital forensics, biometric authentication, and content moderation platforms.

Despite significant advances in in-domain detection, the primary failure mode of contemporary deepfake detectors remains their vulnerability to \emph{cross-dataset generalization collapse} \citep{zi2020wilddeepfake, li2020celeb}. Detectors optimized solely on benchmark training distributions such as FaceForensics++ (FF++) \citep{rossler2019faceforensics} routinely fail when evaluated against novel generation artifacts, alternative face reenactment codecs, unseen demographic distributions, and heavy real-world compression schemes. Prior studies attribute this vulnerability to detectors over-relying on dataset-specific semantic artifacts (e.g., landmark misalignments or distinct facial boundaries) rather than learning intrinsic, generator-agnostic physical inconsistencies.

To combat domain overfitting, recent literature has increasingly adopted large-scale pretrained Vision Transformers (ViTs), notably multimodal foundation models such as CLIP \citep{radford2021learning}. For example, GenD \citep{gend2024} demonstrated that freezing the pretrained transformer backbone and fine-tuning only the LayerNorm parameters (Parameter-Efficient Fine-Tuning, or PEFT) yields strong cross-manipulation transfer. However, single-frame CLIP-based detectors inherently possess two critical limitations:
\begin{enumerate}[leftmargin=*]
    \item \textbf{High-Frequency Residual Filtering:} Standard self-attention layers act predominantly as low-pass semantic filters, obscuring microscopic spatial inconsistencies (such as pixel grid discontinuities and phase errors) produced by upsampling operations.
    \item \textbf{Temporal Blindness:} By evaluating isolated video frames independently ($T=1$), single-frame detectors fail to register inter-frame temporal incoherence, flickering across facial boundaries, and unnatural eyelid/mouth dynamics that persist even in highly compressed videos.
\end{enumerate}

\paragraph{Our Contributions.}
In this work, we present a rigorous, matched-pipeline investigation into pretrained Vision Transformer generalization for video deepfake detection. Under pre-registered statistical evaluation protocols and confirmatory validation on \textbf{FakeAVCeleb\_v1.2} ($N=2{,}998$ videos across 499 authentic identities and 2,499 visual fakes; frozen at tag \texttt{prereg-freeze-h7}), we examine whether architectural additions—namely spatio-temporal attention and high-frequency forensic residual streams—provide measurable benefits over parameter-efficient single-frame tuning, and audit detector reliability under operating point and distribution shifts.

Specifically, this paper delivers the following contributions:
\begin{itemize}[leftmargin=*]
    \item \textbf{Controlled Matched-Pipeline Multi-Seed Benchmark:} Across 12 fully trained models (3 independent random seeds across 4 matched architectures on FaceForensics++ c23) and 1 official foundation checkpoint (\texttt{yermandy/GenD\_CLIP\_L\_14}), we conduct a strictly controlled ablation separating architectural capacity from optimization recipe. We benchmark across in-domain FF++ c23 ($N=840$), Celeb-DF-v2 ($N=518$), DeepFakeDetection ($N=2{,}847$ pure), and the pre-registered confirmatory dataset FakeAVCeleb\_v1.2 ($N=2{,}998$).
    \item \textbf{Recipe vs. Architecture Disentanglement:} Comparing the single-frame LayerNorm baseline (\texttt{smsdt\_s}, $777\text{K}$ parameters) against our in-house reproduction of GenD (\texttt{gend\_ours}, $103\text{K}$ parameters) reveals a $+5.32\%$ AUC gap on Celeb-DF v2 ($0.8826$ vs $0.8294$), $+5.30\%$ on FakeAVCeleb ($0.9437$ vs $0.8907$), and $+5.83\%$ on FF++, demonstrating that training recipe, loss formulation, and sample scheduling drive substantial performance variances previously attributed to architecture.
    \item \textbf{Confirmatory Architectural Equivalence (H1 Supported):} Using paired cluster bootstrap tests ($N_{\text{boot}}=10{,}000$, equivalence boundary $\delta=0.01$, Bonferroni $m=2$) across $991$ clusters on FakeAVCeleb, we confirm that scaling trainable heads by an order of magnitude to spatio-temporal attention (\texttt{smsdt\_st}, $8.5\text{M}$ params; $\Delta\text{AUC} = -0.0045$, $97.5\%\text{ CI }[-0.0096, +0.0006]$) and dual forensic streams (\texttt{smsdt\_full}, $9.1\text{M}$ params; $\Delta\text{AUC} = +0.0012$, $97.5\%\text{ CI }[-0.0050, +0.0073]$) yields equivalence within $\pm 0.01$ $\Delta\text{AUC}$ relative to single-frame PEFT (\texttt{smsdt\_s}), definitively supporting the pre-registered equivalence hypothesis \textbf{H1}.
    \item \textbf{Audit of False Alarm Inflation and GMM Collapse (H4, H5):} We uncover a severe operational failure mode: freezing source thresholds at $\alpha = 5\%$ inflates target false alarms up to $15.7\%$ on Celeb-DF and $32.1\%$ on FakeAVCeleb, and official GenD weights exhibit identical inflation ($15.7\%$ Celeb-DF FPR). Unsupervised 2-component GMM thresholding collapses completely under high fake prevalence ($\pi_{\text{fake}}=0.91$), with $100\%$ degenerate fits (\textbf{H5 falsified}).
    \item \textbf{Split-Conformal Real-Shot Calibration (H7-v2 Supported):} To resolve operational false alarm inflation without full target retraining, we introduce a pre-registered target-real budget transfer benchmark. We prove that split-conformal calibration on a modest budget of $n=20$ pristine target identities guarantees target false alarm control ($\le 4.72\%$ empirical FPR across all 13 models) while preserving $83.4\%\text{--}96.8\%$ of oracle recall, whereas naive quantile thresholding violates the budget ($8.5\%\text{--}9.0\%$ FPR).
    \item \textbf{Unsupervised Test-Time Stabilization:} We evaluate test-time adaptation (TTA) strategies and show that standard online entropy minimization suffers from operational collapse on imbalanced streams, whereas batch-episodic adaptation and a multi-clip consistency objective preserve ranking stability.
\end{itemize}

